\documentclass[11pt]{article}
\usepackage{fontspec}
\usepackage[french,english]{babel}
\usepackage[leqno]{amsmath}
\usepackage{amssymb}
\usepackage{amsthm}
\usepackage{mathrsfs}
\usepackage{tikz-cd}
\usepackage{enumitem}

\newcommand{\printedpage}[1]{\marginpar{\footnotesize p.~#1}}
\newcommand{\msp}[1]{\marginpar{\footnotesize [#1]}}
\newcommand{\cercle}[1]{\tikz[baseline=(c.base)]\node[draw,circle,inner sep=1pt](c){#1};}
\newcommand{\Fibu}{\underline{\mathrm{Fib}}}
\newcommand{\Homu}{\underline{\mathrm{Hom}}}
\newcommand{\Ho}{\mathrm{Ho}}
\newcommand{\ho}{\mathrm{ho}}
\newcommand{\lc}{\mathrm{lc}}

\begin{document}

\printedpage{155}
% head: A. Grothendieck : Dérivateurs, Ch. XII, Édition des Universités de Montpellier II et Paris VII
% foot: Version 30 mars 2003

\msp{page 240}
ersion), and $p_f \in W$. % beginning of a word continued from the preceding page (outside this range), as printed

\section*{7\quad Return to the axiomatics of the $W$-homotopic fiber.}

I come back to the reflection in n$^\circ$ 4. In the end, for
\[
(\mathcal{M}, W \subseteq \mathrm{Fl}(\mathcal{M}))
\]
given with $W$ weakly saturated, it appears that an excellent theory of $W$-homotopic fibers, satisfying the conditions \cercle{a}\cercle{b}\cercle{c} of n$^\circ$ 4 (pages 215-218), exists as soon as the following innocent-looking axiom is satisfied (in addition to the existence of the final object and fiber products (\footnote[128]{and even that! it seems to me that it suffices to assume that the $X \longrightarrow Y$ which are $W$-fibrant are ``squarable''.})).

\textbf{Left factorization:} (\footnote[129]{Under these conditions, $\mathcal{M}$ equipped with $(W, F_W)$ is a category with cofibrations of K.S. \textsc{Brown-Anderson}, and ``everything is O.K.'', but one would have to check a) b) c) above (p. 215, 218), I am not sure I have done it.},\footnote[130]{\textbf{NB} The left factorization axiom is sufficient for the existence of a good theory of homotopy fibers (it seems) only in the case where every $X$ in $\mathcal{M}$ is $W$-fibrant, i.e. is in $\Fibu_W e$.}). Every arrow $f : X \longrightarrow Y$ of $\mathcal{M}$ factors as
\[
X \xrightarrow{i_f} \bar{X} \xrightarrow{p_f} Y,
\]
with
\[
i_f \in W, \qquad p_f \in \mathrm{Fib}_W.
\]

The dual axiom gives rise to a good theory of $W$-homotopic cofibers.

\textbf{Right factorization:} Every arrow $f : X \longrightarrow Y$ in $\mathcal{M}$ factors as
\[
X \xrightarrow{j_f} \tilde{Y} \xrightarrow{q_f} Y,
\]
with
\[
q_f \in W, \qquad j_f \in \mathrm{Cofib}_W.
\]

\msp{page 241}
I would now like to develop somewhat the formalism of $W$-homotopic fibers, under the assumption (say) that $W$ satisfies the left factorization axiom.

Let us consider a cartesian diagram
\[
\begin{tikzcd}[ampersand replacement=\&]
X_1 \arrow[d, "p_1"'] \& Y \arrow[l, "q_2"'] \arrow[d, "q_1"] \\
S \& X_2 \arrow[l, "p_2"]
\end{tikzcd}
\]
\printedpage{156}
% head: A. Grothendieck : Dérivateurs, Ch. XII, Édition des Universités de Montpellier II et Paris VII
% foot: Version 30 mars 2003
in $\mathcal{M}$. If $p_1$ is in $\mathrm{Fib}_W$, then $Y$ over $X_2$ can be interpreted as the $W$-homotopic fiber of $X_1/S$ at $X_2/S$. More precisely, this is so for the class $\ho^{\lc}_{W,X_2}(Y)$,
\[
\ho^{\lc}_{W,X_2}(Y) \simeq p_2^*(\ho^{\lc}_{W,S}(X_1/S)). \tag*{$(*)$}
\]
But as an object over $X_1$, $Y$ can also be interpreted as a $W$-homotopic fiber of $X_2/S$ at $X_1/S$:
\[
\ho^{\lc}_{W,X_1}(Y/X_1) \simeq p_1^*(\ho^{\lc}_{W,S}(X_2/S)). \tag*{$(**)$}
\]
To tell the truth, one would have to prove it. For this, let us consider a factorization
\[
X_2 \xrightarrow{i_2} \bar{X}_2 \xrightarrow{\bar{p}_2} S \qquad i_2 \in W\ p_2 \in \mathrm{Fib}_W, % [sic] p_2 \in \mathrm{Fib}_W -> \bar{p}_2 \in \mathrm{Fib}_W, the factorisation just written is through \bar{p}_2
\]
so that we find a diagram of cartesian squares
\[
\begin{tikzcd}[ampersand replacement=\&]
X_1 \arrow[d] \& \bar{Y} \arrow[l, "\bar{q}_2"'] \arrow[d] \& Y \arrow[l, "j_2"', "W"] \arrow[d] \\
S \& \bar{X}_2 \arrow[l, "\bar{q}_2"'] \& X_2. \arrow[l, "i_2"', "W"] % [sic] \bar{q}_2 (bottom left arrow) -> \bar{p}_2, the arrow \bar{X}_2 \to S is \bar{p}_2 of the factorisation above
\end{tikzcd}
\]

\msp{page 242}
By definition of $p_1^*$, we have
\[
p_1^*(\ho^{\lc}_{W,S}(X_2/S)) \simeq \ho^{\lc}_{W,X_1}(\bar{Y}/X_1),
\]
on the other hand, since $X_1 \longrightarrow S$ is in $\mathrm{Fib}_W$, and $i_2 \in W$, we have $j_2 \in W$, hence
\[
\ho^{\lc}_{W,X_1}(Y/X_1) \xrightarrow{\sim} \ho^{\lc}_{W,X_1}(\bar{Y}/X_1),
\]
whence formula $(**)$.

Thus, the object $Y$ has two different homotopical interpretations, depending on whether one regards it as being over $X_1$, or over $X_2$. One would then like an interpretation which encompasses both at once. We are going to consider the diagram
\[
\begin{tikzcd}[ampersand replacement=\&]
X_1 \arrow[d] \& \\
S \& X_2 \arrow[l]
\end{tikzcd}
\]
of type $\Psi^\circ$, or
\[
\Psi \quad = \quad
\begin{tikzcd}[ampersand replacement=\&]
1 \& \\
0 \arrow[u] \arrow[r] \& 2
\end{tikzcd}
\qquad \text{or} \qquad
\begin{tikzcd}[ampersand replacement=\&, row sep=small, column sep=small]
 \& 1 \\
0 \arrow[ur] \arrow[dr] \& \\
 \& 2
\end{tikzcd}
\quad ,
\]
% the last "2" of the first diagram is printed overlapping the page number 156
\printedpage{157}
% head: A. Grothendieck : Dérivateurs, Ch. XII, Édition des Universités de Montpellier II et Paris VII
% foot: Version 30 mars 2003
from which $Y$ is obtained by fiber product, by completing this diagram to a diagram of type $\tilde{\Psi}^\circ$, where
\[
\tilde{\Psi} =
\begin{tikzcd}[ampersand replacement=\&, row sep=small, column sep=small]
 \& 1 \arrow[dr] \& \\
0 \arrow[ur] \arrow[dr] \& \& 3 \\
 \& 2 \arrow[ur] \&
\end{tikzcd}
\qquad (\simeq \boldsymbol{\Delta}^1 \times \boldsymbol{\Delta}^1).
\]

\msp{page 243}
For every category $I$, let
\[
\Ho_{\mathcal{M},W}(I) \text{ \ or \ } \Ho_W(I) \overset{\text{def}}{=} \mathcal{M}(I)W(I)^{-1},
\]
where
\[
\left\{
\begin{array}{rcl}
\mathcal{M}(I) & = & \Homu(I^\circ, \mathcal{M}) \\[1ex]
W(I) & \subseteq & \mathrm{Fl}(\mathcal{M}(I)) \text{ consists of the } u : F \longrightarrow G \text{ such that} \\[1ex]
 & & u_i : F(i) \longrightarrow G(i) \text{ in } W \text{ for every } i \text{ in } \mathrm{Ob}\, I.
\end{array}
\right.
\]

Thus, the diagram $\vcenter{\hbox{\begin{tikzcd}[ampersand replacement=\&, row sep=small, column sep=small]
 \& X_1 \arrow[dl] \\
S \& \\
 \& X_2 \arrow[ul]
\end{tikzcd}}}$ defines an object $\xi$ of $\mathcal{M}(\Psi)$, hence of $\Ho_W(\Psi)$, and the diagram $\vcenter{\hbox{\begin{tikzcd}[ampersand replacement=\&, row sep=small, column sep=small]
 \& X_1 \arrow[dl] \& \\
S \& \& Y \arrow[ul] \arrow[dl] \\
 \& X_2 \arrow[ul] \&
\end{tikzcd}}}$ can likewise be interpreted as an object $\tilde{\xi}$ of $\Ho_W(\tilde{\Psi})$, whose restriction to $\Psi$ is $\xi$.

This being laid down, I would like to define a functor
\[
i_* : \Ho_W(\Psi) \longrightarrow \Ho_W(\tilde{\Psi})
\]
(where
\[
i : \Psi \longrightarrow \tilde{\Psi}
\]
is the inclusion, which is an open immersion), which must correspond, in the case of an object $\vcenter{\hbox{\begin{tikzcd}[ampersand replacement=\&, row sep=small, column sep=small]
 \& X_1 \arrow[dl] \\
S \& \\
 \& X_2 \arrow[ul]
\end{tikzcd}}}$ of $\Ho_W(\Psi)$ with $p_1 \in \mathrm{Fib}_W$ \texttt{[}$:$ $X_1 \longrightarrow S$\texttt{]}, to the fiber product operation (\footnote[131]{\textbf{NB} This construction, which is done here rather laboriously, becomes immediate, and its main properties too, if one admits that the category with fibrations $(\mathcal{M}, W, \mathrm{Fib}_W)$ defines a \textit{derivator} on the domain of finite ordered sets, which can indeed be proved quite simply (with no more trouble than I am taking here).}).

\printedpage{158}
% head: A. Grothendieck : Dérivateurs, Ch. XII, Édition des Universités de Montpellier II et Paris VII
% foot: Version 30 mars 2003

\msp{page 244}
The construction of this functor $i_*$ should be essentially sorital, with what we have at our disposal. We shall first define, for $S$ \textit{fixed,} a canonical functor
\[
\Ho^{\lc}_W(S) \times \Ho^{\lc}_W(S) \longrightarrow \Ho_{W,S}(\Psi), \tag*{$(*)$}
\]
where the second member is defined, by analogy with the $\Ho_W(I)$, by
\[
\Ho_{W,S}(I) \overset{\text{def}}{=} ((\mathcal{M}/S)(I))(W_S(I))^{-1}
\]
- i.e. it is the same definition, except that $(\mathcal{M}, W)$ is replaced by $(\mathcal{M}/S, W_S)$. As for $\Phi = \vcenter{\hbox{\begin{tikzcd}[ampersand replacement=\&, row sep=small, column sep=small]
1 \arrow[dr] \& \\
 \& 3 \\
2 \arrow[ur] \&
\end{tikzcd}}}$ , we regard it as a closed subcategory of $\tilde{\Psi}$, which is obtained from it by adjoining an initial object 0 (just as $\tilde{\Psi}$ is obtained from $\Psi$ by adjoining a final object 3). On the other hand, we have a tautological functor
\[
\Ho_{W,S}(\Phi) \longrightarrow \Ho_W(\tilde{\Psi}), \tag*{$(*')$}
\]
which on objects consists in completing a diagram $\vcenter{\hbox{\begin{tikzcd}[ampersand replacement=\&, row sep=small, column sep=small] % [sic] les objet -> les objets, plural agreement
X_1 \& \\
 \& Y \arrow[ul] \arrow[dl] \\
X_2 \&
\end{tikzcd}}}$ of $\mathcal{M}/S$ to the square diagram $\vcenter{\hbox{\begin{tikzcd}[ampersand replacement=\&, row sep=small, column sep=small]
 \& X_1 \arrow[dl] \& \\
S \& \& Y \arrow[ul] \arrow[dl] \\
 \& X_2 \arrow[ul] \&
\end{tikzcd}}}$ of $\mathcal{M}$.

\msp{page 245}
The composite of $(*)$ and $(*')$ will therefore give a functor
\[
\Ho^{\lc}_W(S) \times \Ho^{\lc}_W(S) \longrightarrow \Ho_W(\tilde{\Psi}), \tag*{$(**)$}
\]
already very close to the functor $i_*$ to be constructed. To construct $(*)$, we interpret (up to equivalence of categories) the first member as
\[
(\Fibu_W S \times \Fibu_W S)(W_S \times W_S)^{-1},
\]
so we have to construct a functor
\[
\Fibu_W S \times \Fibu_W S \longrightarrow \Ho_{W,S}(\Phi),
\]
which transforms arrows in $W_S \times W_S$ into isomorphisms. We shall take the functor
\[
(X_1, X_2) \ \longmapsto \
\left(
\begin{tikzcd}[ampersand replacement=\&]
X_1 \& X_1 \times_S X_2 \arrow[l] \arrow[d] \\
 \& X_2
\end{tikzcd}
\right),
\]
\printedpage{159}
% head: A. Grothendieck : Dérivateurs, Ch. XII, Édition des Universités de Montpellier II et Paris VII
% foot: Version 30 mars 2003
and we only have to check that if
\[
X_1 \xrightarrow{\alpha_1} X_1', \qquad X_2 \xrightarrow{\alpha_2} X_2'
\]
are two arrows in $\Fibu_W S$, whence
\[
X_1 \times_S X_2 \xrightarrow{\alpha = \alpha_1 \times_S \alpha_2} X_1' \times_S X_2',
\]
then
\[
\alpha_1 \in W,\ \alpha_2 \in W \qquad \Longrightarrow \qquad \alpha \in W,
\]
which is immediate from the factorization
\[
X_1 \times_S X_2 \quad \underbrace{\longrightarrow}_{\substack{\in W \\ \text{since } X_2 \in \Fibu_W S}} \quad X_1' \times_S X_2 \quad \underbrace{\longrightarrow}_{\substack{\in W \\ \text{since } X_1' \in \Fibu_W S}} \quad X_1' \times_S X_2'.
\]

\msp{page 246}
We have proved the

\textbf{Proposition 18.} \textit{Let $S$ be in $\mathcal{M}$. There exists, up to unique isomorphism, a functor and only one}
\[
\Ho^{\lc}_W(S) \times \Ho^{\lc}_W(S) \xrightarrow{\pi} \Ho_{W,S}(\Phi)
\]
\textit{giving rise to a square commutative up to isomorphism}
\[
\begin{tikzcd}[ampersand replacement=\&, column sep=large]
 \& {\scriptstyle (X_1,X_2)\ \longmapsto\ \begin{array}{@{}c@{}c@{}c@{}} \scriptstyle X_1 & & \\ & \nwarrow & \\ & & \scriptstyle X_1 \times_S X_2 \\ & \swarrow & \\ \scriptstyle X_2 & & \end{array}} \& \\[-3ex]
\Fibu_W S \times \Fibu_W S \arrow[rr] \arrow[d] \& \& (\mathcal{M}/S)(\Phi) \arrow[d] \\
\Ho^{\lc}_W(S) \times \Ho^{\lc}_W(S) \arrow[rr] \& \& \Ho_{W,S}(\Phi) \\[-4ex]
(= (\mathcal{M}/S)(W_S)^{-1} \times (\mathcal{M}/S)(W_S)^{-1}) \& \& (= \mathcal{M}/S(\Phi)(W_S(\Phi))^{-1}).
\end{tikzcd}
\]
\textit{For $X_1$, $X_2$ in $\mathcal{M}/S$, the diagram}
\[
\pi(X_1, X_2) \overset{\text{def}}{=} \pi(\ho^{\lc}_{W,S}(X_1), \ho^{\lc}_{W,S}(X_2))
\]
\textit{is ``computed'' as follows, up to canonical isomorphism: one chooses}
\[
X_1 \xrightarrow{\alpha_1} \bar{X}_1, \qquad X_2 \xrightarrow{\alpha_2} \bar{X}_2 \qquad \textit{arrows in } \mathcal{M}/S
\]
\textit{with}
\[
\alpha_1, \alpha_2 \in W_S, \qquad \bar{X}_1, \bar{X}_2 \in \mathrm{Ob}\, \Fibu_W S,
\]
\textit{which defines}
\[
\pi(\alpha_1, \alpha_2) : \pi(X_1, X_2) \xrightarrow{\sim} \pi(\bar{X}_1, \bar{X}_2) = \ho_{S,W,\Phi}
\left(
\begin{tikzcd}[ampersand replacement=\&]
\bar{X}_1 \& \bar{X}_1 \times_S \bar{X}_2 \arrow[l] \arrow[d] \\
 \& \bar{X}_2
\end{tikzcd}
\right).
\]

\end{document}
